% ============================================
% Rivian - Električna vozila budućnosti
% Autor: Miloš Marjanović
% Predmet: Uvod u informatiku
% ============================================

\documentclass[12pt,a4paper]{article}

% Paketi za srpski jezik i ćirilicu
\usepackage[T2A]{fontenc}
\usepackage[utf8]{inputenc}
\usepackage[serbian]{babel}

% Grafika i slike
\usepackage{graphicx}
\usepackage{wrapfig}
\usepackage{float}

% Matematika
\usepackage{amsmath}
\usepackage{amssymb}
\usepackage{amsthm}

% Boje
\usepackage{xcolor}

% Tabele
\usepackage{array}
\usepackage{booktabs}

% Linkovi
\usepackage{hyperref}
\hypersetup{
    colorlinks=true,
    linkcolor=teal,
    urlcolor=orange
}

% Margine
\usepackage[margin=2.5cm]{geometry}

% Definisanje boja u stilu Rivian-a
\definecolor{riviangreen}{RGB}{26, 61, 61}
\definecolor{rivianorange}{RGB}{244, 165, 0}

% Definisanje teorema, definicija i lema
\theoremstyle{definition}
\newtheorem{definicija}{Дефиниција}[section]
\newtheorem{teorema}{Теорема}[section]
\newtheorem{lema}{Лема}[section]

% ============================================
% ПОЧЕТАК ДОКУМЕНТА
% ============================================
\begin{document}

% Насловна страна
\begin{titlepage}
    \centering
    \vspace*{2cm}
    
    {\Huge\bfseries\color{riviangreen} Ривијан}\\[0.5cm]
    {\LARGE Електрична возила будућности}\\[2cm]
    
    {\Large\textit{Семинарски рад}}\\[1cm]
    {\large Предмет: Увод у информатику}\\[2cm]
    
    \begin{tabular}{rl}
        \textbf{Аутор:} & Милош Марјановић \\
        \textbf{Број индекса:} & МИ25160 \\
        \textbf{Датум:} & \today
    \end{tabular}
    
    \vfill
    
    {\large Факултет информатике}\\
    {\large 2025. година}
\end{titlepage}

% Садржај
\tableofcontents
\newpage

% ============================================
\section{Увод}
% ============================================

\textbf{Ривијан} (енг. \textit{Rivian Automotive, Inc.}) је америчка компанија за производњу електричних возила основана \textcolor{rivianorange}{2009. године}. Компанија се фокусира на израду возила дизајнираних за авантуру и истраживање природе.

У овом раду ћемо анализирати технологију електричних возила, математичке принципе који стоје иза њих, као и предности у односу на традиционална возила са мотором са унутрашњим сагоревањем.

\vspace{0.5cm}

\begin{wrapfigure}{r}{0.5\textwidth}
    \centering
    \includegraphics[width=0.48\textwidth]{slike/r1t.jpg}
    \caption{Ривијан R1T}
    \label{fig:r1t}
\end{wrapfigure}

Ривијан је постао познат по својим иновативним моделима:
\begin{itemize}
    \item \textbf{R1T} -- први масовно произведен електрични пикап
    \item \textbf{R1S} -- електрични СУВ са 7 седишта
    \item \textbf{ЕДВ} -- доставно возило у партнерству са Амазоном
\end{itemize}

Мисија компаније је да створи одрживи транспорт који инспирише људе да истражују свет. Компанија верује да је могуће правити возила која су истовремено еколошки одговорна и узбудљива за вожњу.

\vspace{0.5cm}
\clearpage

% ============================================
\section{Како функционишу електрична возила?}
% ============================================

\subsection{Енергија и снага}

Основни физички принципи електричних возила заснивају се на електромагнетизму и претварању енергије. 

\begin{definicija}[Електрична енергија]
Електрична енергија батерије израчунава се као производ снаге и времена:
\begin{equation}
    E = P \cdot t
\end{equation}
где је $E$ енергија у киловат-сатима (kWh), $P$ снага у киловатима (kW), а $t$ време у сатима.
\end{definicija}

За Ривијан R1T са батеријом од 135 kWh и потрошњом од приближно 0.27 kWh/km, теоретски домет се може израчунати као:

\begin{equation}
    D = \frac{E_{bat}}{E_{cons}} = \frac{135 \text{ kWh}}{0.27 \text{ kWh/km}} = 500 \text{ km}
\end{equation}

\subsection{Обртни момент електромотора}

\begin{teorema}[Обртни момент]
Обртни момент електромотора је максималан од нултих обртаја, што даје тренутну акцелерацију:
\begin{equation}
    \tau = k \cdot I \cdot \Phi
\end{equation}
где је $\tau$ обртни момент, $k$ константа мотора, $I$ јачина струје, а $\Phi$ магнетни флукс.
\end{teorema}

\begin{lema}[Ефикасност]
Ефикасност електромотора износи типично 85-95\%, што је значајно више од мотора са унутрашњим сагоревањем (20-30\%):
\begin{equation}
    \eta_{EM} = \frac{P_{izlaz}}{P_{ulaz}} \approx 0.90
\end{equation}
\end{lema}

% ============================================
\section{Спецификације возила}
% ============================================

У следећој табели приказане су техничке спецификације Ривијан модела:

\begin{table}[H]
\centering
\caption{Поређење спецификација Ривијан модела}
\label{tab:specs}
\begin{tabular}{|l|c|c|c|}
\hline
\textbf{Спецификација} & \textbf{R1T} & \textbf{R1S} & \textbf{ЕДВ 700} \\
\hline
Тип возила & Пикап & СУВ & Комби \\
\hline
Снага (KS) & \textcolor{rivianorange}{835} & \textcolor{rivianorange}{835} & $\sim$400 \\
\hline
Обртни момент (Nm) & 1.231 & 1.231 & $\sim$600 \\
\hline
0-100 km/h (s) & \textbf{3.0} & \textbf{3.0} & N/A \\
\hline
Домет (km) & до 505 & до 505 & до 240 \\
\hline
Батерија (kWh) & 135 & 135 & $\sim$100 \\
\hline
Број седишта & 5 & 7 & 2 \\
\hline
\end{tabular}
\end{table}

% ============================================
\section{Предности електричних возила}
% ============================================

Електрична возила имају бројне предности:

\begin{enumerate}
    \item \textbf{Еколошке предности}
    \begin{itemize}
        \item Нула директних емисија $\text{CO}_2$
        \item Мањи укупни угљенични отисак
        \item Тише функционисање
    \end{itemize}
    
    \item \textcolor{riviangreen}{\textbf{Економске предности}}
    \begin{itemize}
        \item Нижи трошкови горива (електрична енергија је јефтинија)
        \item Мање покретних делова = мање одржавања
        \item Потенцијални порески подстицаји
    \end{itemize}
    
    \item \textit{Перформансне предности}
    \begin{itemize}
        \item Тренутни обртни момент
        \item Нижи центар тежишта
        \item Напредне технологије (аутономна вожња)
    \end{itemize}
\end{enumerate}

% ============================================
\section{Закључак}
% ============================================

Ривијан представља \textcolor{riviangreen}{\textbf{будућност аутомобилске индустрије}}. Комбинацијом иновативног дизајна, напредне технологије и посвећености одрживости, компанија је успела да створи возила која нису само еколошки прихватљива, већ и \textit{узбудљива за вожњу}.

Математички принципи иза електричних возила, као што је Ајнштајнова формула $E = mc^2$ за еквиваленцију масе и енергије, показују колико је енергије заправо садржано у батеријама.

Будућност транспорта је \textcolor{rivianorange}{\textbf{електрична}}, а Ривијан је на челу те револуције.


\end{document}
